%! Projections onto Subspaces — Unit 4, Lesson 4.2 — Lesson Plan
\documentclass[10pt]{article}
\usepackage{linalg-article}
\usepackage{linalg-boxes}
\usepackage{graphicx}   % -article does NOT load graphicx; needed for thumbnails/screenshots
\graphicspath{{images/}}
\newcommand{\TallMath}[1]{$\displaystyle #1\rule[-1.4em]{0pt}{3.2em}$}
\newcommand{\vv}{\mathbf{v}}
\newcommand{\ww}{\mathbf{w}}
\newcommand{\xx}{\mathbf{x}}
\newcommand{\bb}{\mathbf{b}}
\newcommand{\pp}{\mathbf{p}}
\newcommand{\ee}{\mathbf{e}}
\newcommand{\avec}{\mathbf{a}}
\newcommand{\zero}{\mathbf{0}}
\newcommand{\T}{\mathsf{T}}
\newcommand{\UnitNumberName}{Unit 4: Orthogonality \quad}
\newcommand{\LessonNumberName}{Lesson 4.2: Projections onto Subspaces}

\begin{document}
\begin{center}
    {\Large\bfseries \CourseName: \SchoolYear \\
    {\small\bfseries \UnitNumberName \LessonNumberName}}
\end{center}

\begin{tcolorbox}[colback=blush, colframe=burgundy, boxrule=0.9pt, arc=2mm,
    left=2mm, right=2mm, top=1.5mm, bottom=1.5mm]
    \textbf{Primary Objective:} Students can find the \emph{projection} $\pp$ of a vector $\bb$ onto a
    line or a subspace --- the closest point --- by forcing the \emph{error} $\ee=\bb-\pp$ to be
    perpendicular to the space. They can derive and use $\hat{x}=\frac{\avec\cdot\bb}{\avec\cdot\avec}$
    for a line and the \textbf{normal equations} $A^{\T}A\hat{\xx}=A^{\T}\bb$ for a subspace, and can
    explain \emph{why} ``closest'' means ``error perpendicular'' --- the error lands in the orthogonal
    complement of the space (Lesson~4.1).
\end{tcolorbox}

\renewcommand{\arraystretch}{1.4}
\begin{skillbox}[Priority Ideas \& Skills]{goldbox}
\begin{tabularx}{\linewidth}{@{}X|X@{}}
\textbf{Priority Ideas \& Skills} & \textbf{Key Understandings (the ``why'')} \\[2pt]
\hline
\vspace{-6pt}
\begin{itemize}[leftmargin=1.1em, nosep, after=\vspace{2pt}]
  \item Find the closest point $\pp$ on a \textbf{line} to a target $\bb$: $\hat{x}=\frac{\avec\cdot\bb}{\avec\cdot\avec}$, $\pp=\hat{x}\avec$.
  \item Read $\pp=P\bb$ off a \textbf{projection matrix} $P=\frac{\avec\avec^{\T}}{\avec^{\T}\avec}$.
  \item Project onto a \textbf{subspace} $C(A)$ with the \textbf{normal equations} $A^{\T}A\hat{\xx}=A^{\T}\bb$.
  \item Locate the \textbf{error} $\ee=\bb-\pp$ and check $\ee\perp$ the space.
\end{itemize}
&
\vspace{-6pt}
\begin{itemize}[leftmargin=1.1em, nosep, after=\vspace{2pt}]
  \item ``Closest'' $=$ ``error perpendicular'': drop a perpendicular to the space (Lesson~4.0).
  \item $\avec\cdot\ee=0$ is one dot-product equation; $A^{\T}\ee=\zero$ is one per column.
  \item The error lands in the \emph{orthogonal complement} $N(A^{\T})$ of the column space (Lesson~4.1).
  \item Projecting an already-projected point changes nothing: $\pp$ is already in the space.
\end{itemize}
\end{tabularx}
\end{skillbox}

\begin{skillbox}[{Vocabulary, Concepts, \& Theorems}]{greenbox}
\renewcommand{\arraystretch}{1.3}
\begin{tabularx}{\linewidth}{@{}l X@{}}
\textbf{Projection $\pp$} & the closest point to $\bb$ that lies \emph{in} the line or subspace. \\
\textbf{Error (residual) $\ee$} & $\ee=\bb-\pp$; it is perpendicular to the whole space ($\avec\cdot\ee=0$, or $A^{\T}\ee=\zero$). \\
\textbf{Projection onto a line} & $\hat{x}=\dfrac{\avec\cdot\bb}{\avec\cdot\avec}$ and $\pp=\hat{x}\,\avec$. \\
\textbf{Projection matrix $P$} & $\pp=P\bb$; a line gives $P=\dfrac{\avec\avec^{\T}}{\avec^{\T}\avec}$, a subspace $P=A(A^{\T}A)^{-1}A^{\T}$. \\
\textbf{Normal equations} & $A^{\T}A\hat{\xx}=A^{\T}\bb$: solve for the weights $\hat{\xx}$, then $\pp=A\hat{\xx}$. \\
\end{tabularx}
\end{skillbox}

\begin{fixedskillbox}[Activate Prior Knowledge \& Spiral Review]{blush}
    The Warm-Up rehearses the three skills this lesson builds on:
    \textbf{(1)} the \emph{perpendicular test} $\vv\cdot\ww=0$ (Lessons~1.2, 4.0);
    \textbf{(2)} computing the two dot products $\avec\cdot\bb$ and $\avec\cdot\avec$ (the pieces of
    $\hat{x}$); and \textbf{(3)} \emph{solving a $2\times2$ system} (the shape of the normal equations).
    Debrief item~3 aloud: emphasize that projecting onto a subspace ends in an ordinary linear system ---
    the machinery is all from Units~1--2.
\end{fixedskillbox}

\begin{skillbox}[Hook]{blush}
    A drone hovers at a point $\bb$ off to the side of a straight road (a \emph{line} through the
    origin). Ask: \textbf{``Which point of the road is closest to the drone, and how do you know it is
    the closest?''} Students reason to the \emph{foot of the perpendicular} --- the drone's shadow
    straight down onto the road. Name the payoff: the closest point is $\pp$ (the \textbf{projection}),
    and the segment from $\pp$ to $\bb$ is the \textbf{error} $\ee=\bb-\pp$, which is \emph{perpendicular}
    to the road. Then pose the lesson question: \textbf{``If `closest' always means `error perpendicular,'
    can we turn that one right angle into an equation and just solve for $\pp$?''}
\end{skillbox}

\begin{skillbox}[Lesson]{blush}
\begin{multicols}{2}
\textbf{1. Closest $=$ perpendicular error.} A target $\bb$ sits off a line through $\avec$. The closest
point $\pp$ on the line is the one where the leftover $\ee=\bb-\pp$ points \emph{straight at} the line ---
$\ee\perp\avec$. \textbf{Ask:} why perpendicular? Any other point on the line is farther (hypotenuse
beats leg). This is ``drop a perpendicular'' from Lesson~4.0, now made exact.

\textbf{2. Projection onto a line.} Write $\pp=\hat{x}\avec$ (some multiple of $\avec$). Force the right
angle: $\avec\cdot(\bb-\hat{x}\avec)=0$. Solve for the one unknown:
$\hat{x}=\dfrac{\avec\cdot\bb}{\avec\cdot\avec}$, so $\pp=\dfrac{\avec\cdot\bb}{\avec\cdot\avec}\,\avec$.
\emph{One dot-product equation, one unknown.} Numbers: $\avec=(1,2)$, $\bb=(4,3)$ give $\hat{x}=2$,
$\pp=(2,4)$, $\ee=(2,-1)$, and $\avec\cdot\ee=0$.

\textbf{3. The projection matrix.} Since $\pp=P\bb$, a line gives $P=\frac{\avec\avec^{\T}}{\avec^{\T}\avec}$
--- one matrix that projects \emph{any} $\bb$. \textbf{Ask:} what is $P\pp$? Still $\pp$ --- projecting
an already-on-the-line point changes nothing.

\columnbreak

\textbf{4. Projection onto a subspace.} Now the target's ``road'' is a plane $=C(A)$ (columns of $A$).
We want $\pp=A\hat{\xx}$ closest to $\bb$, so the error is perpendicular to \emph{every} column:
$A^{\T}(\bb-A\hat{\xx})=\zero$. Rearranged, that is the \textbf{normal equations}
\[ A^{\T}A\,\hat{\xx}=A^{\T}\bb. \]
Solve this ordinary system for $\hat{\xx}$, then $\pp=A\hat{\xx}$. \emph{One dot-product equation per
column} --- the line case, repeated.

\textbf{5. Where the error goes.} $A^{\T}\ee=\zero$ says $\ee$ is in the \textbf{left nullspace}
$N(A^{\T})$ --- the \emph{orthogonal complement} of the column space (Lesson~4.1). \textbf{Punchline:}
projection splits $\bb$ into a piece \emph{in} the space ($\pp$) plus a piece \emph{perpendicular} to it
($\ee$). Forcing $\ee$ into the complement (dot $=0$) is the whole method --- and it is exactly the
setup for least squares in 4.3.
\end{multicols}
\end{skillbox}

\begin{skillbox}[Active Monitoring]{redbox}
Circulate during the notes and Tier work. Watch for and correct:
\begin{itemize}[leftmargin=1.2em, nosep]
  \item flipping the fraction to $\frac{\avec\cdot\avec}{\avec\cdot\bb}$ --- the unknown multiplies $\avec$, so $\avec\cdot\avec$ is on the bottom;
  \item projecting onto $\bb$ instead of onto $\avec$ (the space is spanned by $\avec$, not by the target);
  \item in the subspace case, forgetting to multiply through by $A^{\T}$ --- $A\hat{\xx}=\bb$ usually has \emph{no} solution, which is the whole reason we project.
\end{itemize}
Cold-call prompts: ``Why is the closest point the one with a perpendicular error?'' ``What does
$\avec\cdot\ee=0$ let you solve for?'' ``Which subspace does the error $\ee$ live in, and why?''
\end{skillbox}

\begin{skillbox}[Group Work \& Differentiation]{redbox}
\textbf{Drop the Perpendicular} activity (tiered; mirrors the notes).
\begin{multicols}{3}
\textbf{Tier R --- Remediate.} Project a target onto a line in $\mathbb{R}^2$: compute $\hat{x}$, $\pp$,
and $\ee$, and check $\avec\cdot\ee=0$.

\columnbreak
\textbf{Tier A --- Approaching.} Project onto a line in $\mathbb{R}^3$; build the projection matrix
$P=\frac{\avec\avec^{\T}}{\avec^{\T}\avec}$ and confirm $P\bb=\pp$.

\columnbreak
\textbf{Tier E --- Extension.} Project onto a \emph{plane} $C(A)$: set up and solve the normal equations
$A^{\T}A\hat{\xx}=A^{\T}\bb$, find $\pp$, and verify $\ee\perp$ both columns.
\end{multicols}
\end{skillbox}

\begin{skillbox}[Individual Work \& Assessment]{redbox}
\textbf{Exit Ticket} (independent, no notes): project a target onto a line and give $\hat{x}$ and $\pp$;
a \textbf{conceptual/justification check} --- say what the error $\ee$ measures and why it is
perpendicular; and name the equation you solve to project onto a plane $C(A)$ and justify it from
``error $\perp$ every column.'' Sort responses into ``sets up $\avec\cdot\ee=0$ / flips the fraction /
projects onto $\bb$'' to decide whether 4.3 opens with a quick re-anchor on the perpendicular-error idea.
\end{skillbox}

\begin{skillbox}[Reinforcement \& Extension]{goldbox}
\textbf{Homework:} project onto a line in $\mathbb{R}^2$ and $\mathbb{R}^3$ (with the projection matrix);
set up and solve the normal equations to project onto a plane; and explain in words why the error is
perpendicular. \textbf{Extension:} verify ``project twice $=$ project once'' ($P\pp=\pp$) and check the
error lands in $N(A^{\T})$, tying back to 4.1. \textbf{Preview:} Lesson~4.3, \emph{Least Squares} ---
when $A\xx=\bb$ has no solution, the projection $\pp$ gives the \emph{best-fit} answer, and the normal
equations become the tool for fitting a line to real data.
\end{skillbox}
\end{document}
